\chapter{Conclusion and Future Work}
\label{ch:conclusion}

\section{Introduction}

This final chapter draws the thesis to a close. It revisits the objectives set in
Chapter~1 and judges how far each was met, summarises the main findings, sets out
the limitations of the work honestly, and proposes directions for future
development. It ends with a few concluding remarks.

\section{Summary of the Work}

This project set out to design and build a practical platform that lets a user
understand, question, visualise and map a private collection of documents through
one web application. That aim has been met. The platform, Synapse, ingests a mixed
collection of documents through a resumable pipeline that handles born-digital,
scanned and structured files; answers questions through a corrective, multi-agent
chat that grounds every claim in a citation and abstains when the evidence is
insufficient; generates faithful artefacts through an agent that binds data in
code; maps a library as an interactive knowledge graph; and supports teams who
share libraries and collaborate in real time.

Each of the objectives from Chapter~1 was achieved. The ingestion pipeline
processes mixed documents reliably (Objective~1, Chapter~5). The grounded chat
answers from the user's library, cites its sources and abstains when it should
(Objective~2, Chapters~4 and~7). The agent plans its work, clarifies only when
necessary, and produces charts, documents and images (Objective~3, Chapters~4
and~6). The knowledge-graph module extracts and displays the entities and
relationships of a library (Objective~4, Chapters~4 and~6). Teams share libraries
and chat in real time (Objective~5, Chapter~6). Finally, the platform was
evaluated on an established benchmark and on purpose-built datasets, and its
accuracy, honesty and efficiency were reported (Objective~6, Chapter~7).

\section{Key Findings}

Three findings stand out. First, a disciplined pipeline that combines layout
detection, optical character recognition and a vision-language model can process
scanned and structured documents well enough that retrieval over them is strong,
with only a small gap from born-digital documents. Second, grading evidence before
answering, and allowing the system to abstain, measurably improves honesty: the
abstention mechanism roughly halved the rate of confident mistakes, while
citations remained reliable. Third, keeping the language model's role to describing
artefacts, and binding data to them in code, produces charts and documents whose
figures are faithful to the source, which is what makes the agent useful in
practice.

\section{Limitations}

The work has several limitations, which are stated plainly. The hardest questions,
those that require evidence to be combined across several documents, remain the
weakest area, as they are for current systems generally. Answering structured
documents occasionally fails on questions that require aggregating across many
rows, where exact computation would help. A full, deep answer can take longer than
a user might like on the most demanding questions, because the careful reasoning
that makes answers reliable takes time. The platform depends on external model
services, so its behaviour is tied to theirs, and the graphics-intensive processing
runs on a single server, which bounds how much work can be in progress at once.

\section{Future Work}

These limitations suggest clear directions for future work. Multi-hop reasoning
could be strengthened by iterating retrieval more deliberately, gathering evidence
across several rounds before answering. The handling of structured documents could
be improved by letting the agent compute aggregates exactly over a spreadsheet
rather than relying on retrieval alone. Latency on the hardest questions could be
reduced by caching intermediate results and by running independent reasoning steps
in parallel. The knowledge graph could be used more actively during answering, so
that questions which span a whole corpus draw on its structure. Finally, the
processing tier could be distributed across more than one server to support larger
collections and more concurrent users.

\section{Concluding Remarks}

This thesis has shown that the separate advances in retrieval-augmented generation,
agents, knowledge graphs and document understanding can be brought together into a
single, practical product. Synapse demonstrates that a careful pipeline paired with
a corrective, multi-agent design can deliver trustworthy document intelligence over
a user's own private collection: answers that are grounded and checkable, a system
that knows when to decline, artefacts that are faithful to their sources, and a map
of how a collection's ideas connect. The result is a platform that turns a static
pile of documents into something a person, or a team, can genuinely understand,
question, visualise and explore.
